\documentclass[conference]{IEEEtran}
\usepackage[utf8]{inputenc}
\usepackage[T1]{fontenc}
\usepackage[brazilian]{babel}
\usepackage{cite}
\usepackage{amsmath,amssymb}
\usepackage{graphicx}
\usepackage{booktabs}
\usepackage{url}
\usepackage{textcomp}
\graphicspath{{./}{figures/}}

\begin{document}

\title{Detecção de Anomalias Independente de Classe com Mapas de Calor por Pixel no MVTec AD}

\author{\IEEEauthorblockN{Pedro Fernandes Galindo}
\IEEEauthorblockA{\textit{Centro de Informática (CIn)} \\
\textit{Universidade Federal de Pernambuco}\\
Recife, Brasil \\
pfg@cin.ufpe.br}
\and
\IEEEauthorblockN{Renato Monteiro Correia Medeiros}
\IEEEauthorblockA{\textit{Centro de Informática (CIn)} \\
\textit{Universidade Federal de Pernambuco}\\
Recife, Brasil \\
rmcm@cin.ufpe.br}
}

\maketitle

\begin{abstract}
Avaliamos um detector de anomalias visuais treinado somente com imagens normais das 15 categorias do MVTec AD, sem receber a identidade da categoria. A versão completa acrescenta atenção local mascarada e perturbação de características ao baseline de reconstrução. No teste oficial, obteve 78,46\% de AUROC de imagem, 89,08\% de AUROC de pixel e 78,75\% de AUPRO macro. A diferença entre os escores normais por categoria dificultou o uso de um limiar global.
\end{abstract}

\begin{IEEEkeywords}
detecção de anomalias, independente de classe, localização de anomalias, mapa de calor, MVTec AD
\end{IEEEkeywords}

\section{Introdução}
Em inspeção visual, defeitos são raros e podem assumir formas desconhecidas, dificultando a obtenção de dados rotulados para treinamento. Este projeto considera o cenário não supervisionado: o modelo aprende exclusivamente com imagens normais e deve assinalar desvios em imagens de diferentes categorias. A saída principal será um mapa de calor espacial com um escore de anomalia para cada pixel, permitindo localizar os defeitos, além de resumir cada mapa em um escore por imagem \cite{bergmann2021mvtec}.

Trabalhos recentes mostram que modelos unificados podem operar sobre várias categorias sem receber rótulos de classe e apontam desafios como atalhos de reconstrução e escalas de pontuação incompatíveis entre categorias. A proposta adota essas ideias como referência para desenvolver e avaliar um único detector sem condicionamento por classe, sem pretender reproduzir integralmente os métodos publicados \cite{you2022unified,guo2024cada}.

\section{Objetivos}
\subsection{Objetivo geral}
Desenvolver e avaliar um modelo geral de detecção de anomalias, independente de classe, que aprenda com imagens normais de múltiplas categorias e produza um mapa de calor com um escore de anomalia por pixel.

\subsection{Objetivos específicos}
\begin{itemize}
\item Caracterizar imagens, categorias, partições e máscaras do MVTec AD.
\item Treinar um único modelo com imagens normais de todas as categorias, sem usar rótulos de classe na entrada nem na inferência.
\item Gerar um mapa de calor de escores contínuos por pixel e um escore agregado por imagem.
\item Medir detecção em nível de imagem e localização em nível de pixel, sem usar o teste oficial para escolher o modelo.
\item Analisar falsos alarmes, custo computacional e limitações da abordagem independente de classe.
\end{itemize}

\section{Justificativa}
Treinar um modelo por categoria exige identificar a classe na implantação e manter múltiplos modelos. Um detector único que não dependa dessa informação pode simplificar o uso em cenários nos quais as categorias são desconhecidas, mas precisa lidar com distribuições visuais distintas e produzir escores comparáveis entre elas. UniAD estuda a detecção unificada sem rótulos de classe e CADA aborda o desalinhamento das distribuições de escores nesse cenário \cite{you2022unified,guo2024cada}.

O MVTec AD reúne objetos, texturas e máscaras pixel a pixel, permitindo avaliar tanto se o modelo detecta uma imagem defeituosa quanto se localiza precisamente a região anômala. Assim, o projeto investigará uma saída diretamente útil à inspeção visual: um mapa de calor espacial com escore por pixel, aprendido sem exemplos de defeitos e sem dependência da classe \cite{bergmann2021mvtec}.

\section{Base de Dados}
O MVTec AD \cite{bergmann2021mvtec} reúne 15 categorias, sendo 5 de textura e 10 de objeto. O conjunto de treino contém apenas imagens normais, enquanto o de teste contém imagens normais e defeituosas. No total são 5.354 imagens, valor que confere com o informado no artigo da base. A Tabela~\ref{tab:contagem} apresenta as contagens por categoria e partição.

\begin{table}[htbp]
\caption{Imagens por categoria e partição (conjunto original)}
\label{tab:contagem}
\centering
\footnotesize
\setlength{\tabcolsep}{3pt}
\begin{tabular}{lrrrr}
\toprule
Categoria & Treino normal & Teste normal & Teste defeit. & Total \\
\midrule
bottle     & 209 & 20 & 63  & 292 \\
cable      & 224 & 58 & 92  & 374 \\
capsule    & 219 & 23 & 109 & 351 \\
carpet     & 280 & 28 & 89  & 397 \\
grid       & 264 & 21 & 57  & 342 \\
hazelnut   & 391 & 40 & 70  & 501 \\
leather    & 245 & 32 & 92  & 369 \\
metal\_nut & 220 & 22 & 93  & 335 \\
pill       & 267 & 26 & 141 & 434 \\
screw      & 320 & 41 & 119 & 480 \\
tile       & 230 & 33 & 84  & 347 \\
toothbrush & 60  & 12 & 30  & 102 \\
transistor & 213 & 60 & 40  & 313 \\
wood       & 247 & 19 & 60  & 326 \\
zipper     & 240 & 32 & 119 & 391 \\
\midrule
\textbf{Total} & \textbf{3629} & \textbf{467} & \textbf{1258} & \textbf{5354} \\
\bottomrule
\end{tabular}
\end{table}

Para validação, foram reservados 10\% das imagens normais de treino de cada categoria (arredondado para cima), com semente fixa 42. O teste oficial permanece intocado. A Tabela~\ref{tab:divisao} mostra a divisão resultante, com 3.262 imagens de treino e 367 de validação.

\begin{table}[htbp]
\caption{Divisão treino/validação das imagens normais de treino}
\label{tab:divisao}
\centering
\small
\begin{tabular}{lrr|lrr}
\toprule
Categoria & Treino & Valid. & Categoria & Treino & Valid. \\
\midrule
bottle     & 188 & 21 & pill       & 240 & 27 \\
cable      & 201 & 23 & screw      & 288 & 32 \\
capsule    & 197 & 22 & tile       & 207 & 23 \\
carpet     & 252 & 28 & toothbrush & 54  & 6  \\
grid       & 237 & 27 & transistor & 191 & 22 \\
hazelnut   & 351 & 40 & wood       & 222 & 25 \\
leather    & 220 & 25 & zipper     & 216 & 24 \\
metal\_nut & 198 & 22 & \textbf{Total} & \textbf{3262} & \textbf{367} \\
\bottomrule
\end{tabular}
\end{table}

As 1.258 imagens defeituosas de teste distribuem-se em 73 tipos de defeito.

\section{Análise Exploratória}
A Figura~\ref{fig:contagem} resume a composição de cada categoria. Há desbalanceamento claro: a categoria toothbrush tem apenas 102 imagens, enquanto hazelnut tem 501. A proporção de imagens normais no teste também varia bastante; em transistor, por exemplo, são 60 normais contra 40 defeituosas.

\begin{figure*}[t]
\centering
\includegraphics[width=\textwidth]{imagens_por_categoria.png}
\caption{Imagens do MVTec AD por categoria e partição.}
\label{fig:contagem}
\end{figure*}

A Figura~\ref{fig:categorias} mostra uma imagem normal de treino de cada uma das 15 categorias, cobrindo texturas (carpet, grid, leather, tile, wood) e objetos.

\begin{figure}[htbp]
\centering
\includegraphics[width=\columnwidth]{normais_15_categorias.png}
\caption{Uma imagem normal de treino para cada uma das 15 categorias do MVTec AD.}
\label{fig:categorias}
\end{figure}

A Figura~\ref{fig:exemplos} apresenta exemplos de defeitos com a máscara de ground truth sobreposta, ilustrando tipos e tamanhos distintos de anomalia.

\begin{figure}[htbp]
\centering
\includegraphics[width=\columnwidth]{defeito_bottle_broken_large.png}\\[2pt]
\includegraphics[width=\columnwidth]{defeito_carpet_cut.png}\\[2pt]
\includegraphics[width=\columnwidth]{defeito_transistor_damaged_case.png}
\caption{Imagem normal de treino, imagem de teste com defeito e máscara ground truth sobreposta em vermelho: bottle (broken large), carpet (cut) e transistor (damaged case).}
\label{fig:exemplos}
\end{figure}

\subsection{Resoluções e tamanho dos defeitos}
Todas as imagens de uma mesma categoria têm a mesma resolução, e todas as 1.258 máscaras têm o mesmo tamanho da imagem correspondente (nenhuma divergência). Grid, screw e zipper são em tons de cinza (um canal); as demais são RGB. A Tabela~\ref{tab:perfil} traz a resolução, o modo de cor e a proporção de pixels anômalos por categoria (pixels marcados divididos pelos pixels da imagem).

\begin{table}[htbp]
\caption{Perfil das imagens e proporção de pixels anômalos por categoria (valores em \%)}
\label{tab:perfil}
\centering
\footnotesize
\setlength{\tabcolsep}{3pt}
\begin{tabular}{llcrrrr}
\toprule
Categoria & Resolução & Cor & Média & Mediana & Mín. & $<1\%$ \\
\midrule
bottle & 900$\times$900 & RGB & 7.62 & 5.72 & 0.576 & 3/63 \\
cable & 1024$\times$1024 & RGB & 4.68 & 3.68 & 0.144 & 7/92 \\
capsule & 1000$\times$1000 & RGB & 1.11 & 0.43 & 0.037 & 85/109 \\
carpet & 1024$\times$1024 & RGB & 2.11 & 1.48 & 0.165 & 37/89 \\
grid & 1024$\times$1024 & Cinza & 0.95 & 0.60 & 0.188 & 39/57 \\
hazelnut & 1024$\times$1024 & RGB & 3.35 & 1.62 & 0.233 & 18/70 \\
leather & 1024$\times$1024 & RGB & 0.87 & 0.46 & 0.089 & 70/92 \\
metal\_nut & 700$\times$700 & RGB & 14.49 & 3.67 & 0.436 & 7/93 \\
pill & 800$\times$800 & RGB & 3.98 & 1.03 & 0.050 & 68/141 \\
screw & 1024$\times$1024 & Cinza & 0.34 & 0.29 & 0.099 & 118/119 \\
tile & 840$\times$840 & RGB & 9.80 & 7.93 & 0.842 & 2/84 \\
toothbrush & 1024$\times$1024 & RGB & 2.14 & 1.08 & 0.238 & 15/30 \\
transistor & 1024$\times$1024 & RGB & 11.98 & 2.03 & 0.224 & 6/40 \\
wood & 1024$\times$1024 & RGB & 5.09 & 2.90 & 0.263 & 8/60 \\
zipper & 1024$\times$1024 & Cinza & 2.62 & 2.35 & 0.159 & 16/119 \\
\bottomrule
\end{tabular}
\end{table}

A Figura~\ref{fig:hist} mostra a distribuição da proporção de pixels anômalos, em escala logarítmica. Dos 1.258 defeitos, 499 (cerca de 40\%) ocupam menos de 1\% da imagem, e 17 ocupam menos de 0,1\%. Os defeitos são especialmente pequenos em screw (118 de 119 máscaras abaixo de 1\%, mediana de 0,29\%), capsule (85 de 109) e leather (70 de 92). Em contrapartida, categorias como metal\_nut e transistor têm defeitos grandes, causados sobretudo pelos tipos flip e misplaced, cujas máscaras cobrem cerca de 38\% a 49\% da imagem. Como ponto de partida, será considerada a resolução de $512\times512$, comparando-se o desempenho com a resolução nativa e monitorando especialmente a detecção dos defeitos menores.

\begin{figure}[htbp]
\centering
\includegraphics[width=\columnwidth]{histograma_proporcao_anomala.png}
\caption{Distribuição da proporção de pixels anômalos nas máscaras de ground truth (eixo horizontal em escala logarítmica; linha tracejada em 1\%).}
\label{fig:hist}
\end{figure}

\section{Modelos em Aprendizagem Profunda}
Será implementada uma versão reduzida e própria, inspirada conceitualmente em UniAD \cite{you2022unified}, usando PyTorch \cite{paszke2019pytorch}. A implementação oficial UniAD foi avaliada como referência, mas não será reutilizada diretamente: fixa dependências antigas (PyTorch 1.5.0 e torchvision 0.6.0+cu90), usa operações CUDA explícitas e configura entrada de 224$\times$224 e treinamento de até 1000 épocas, diferentes do ambiente e da resolução candidata deste projeto \cite{you2022uniadcode}. A biblioteca anomalib foi também avaliada; sua versão 2.7.1 não lista UniAD, e o UniNet disponível é outro método com integração MVTec documentada por categoria \cite{anomalib2026,uninetdocs}. Nenhum código dessas duas implementações será copiado. O DataLoader do projeto usa Pillow e NumPy para leitura e conversão de imagens \cite{pillow,numpy2020}.

O modelo mínimo usa Wide-ResNet-50-2, baseado em redes residuais largas \cite{zagoruyko2016wide}, pré-treinado no ImageNet \cite{russakovsky2015imagenet} e congelado, conforme a implementação e os pesos do torchvision \cite{torchvision2026}. A entrada é redimensionada para 512$\times$512, sem recorte. As saídas de \texttt{layer1}, \texttt{layer2} e \texttt{layer3} (256, 512 e 1024 canais) são alinhadas por interpolação bilinear à grade de stride 16 e concatenadas em um tensor de 1792 canais, com grade espacial 32$\times$32. O decodificador próprio usa projeção 1$\times$1 para 128 canais, embedding posicional e 1024 consultas espaciais aprendidas em um Transformer com duas camadas de encoder e duas de decoder, quatro cabeças e feed-forward de dimensão 512. A saída é projetada de volta a 1792 canais e treinada para reconstruir as features congeladas por MSE. Na versão completa, jitter gaussiano escalado pela norma das features é aplicado somente à entrada do decoder no treino; a atenção do encoder e decoder é restrita a vizinhos em uma janela 7$\times$7. O alvo de reconstrução permanece sem perturbação. O mapa por pixel é a norma L2 da discrepância de features por posição, interpolada bilinearmente para a resolução da imagem. Suavização Gaussiana é opcional; a visualização pode ser normalizada, mas métricas e escores usam os valores brutos. O escore por imagem será a média dos 1\% maiores escores locais, decisão feita antes de avaliar o teste por ser mais robusta a um único pixel ruidoso do que o máximo; a alternativa de máximo após suavização Gaussiana também está implementada.

O baseline desativa a máscara local e o \textit{jitter}. A variante completa usa uma vizinhança $7\times7$ e \textit{jitter} de escala 20 e probabilidade 1 no treino. Os dois modelos foram treinados no Apuana com AdamW, taxa inicial $3\times10^{-4}$, batch 4, acumulação de dois passos e semente 42. Os diferentes orçamentos e critérios de seleção limitam a atribuição dos resultados à arquitetura.

\section{Experimentos}
Os splits fixos contêm 3.262 imagens normais de treino e 367 de validação. O baseline chegou à época 111 e foi selecionado pelo menor MSE de validação na época 109. O treino completo parou na época 385 e selecionou a 315, com \textit{early stopping} baseado na média móvel de 25 épocas, melhora mínima de $10^{-4}$ e tolerância de 70 épocas.

Os dois modelos passaram uma vez pelas mesmas 1.725 imagens do teste oficial. Inspeções visuais anteriores de amostras desse conjunto ficaram fora da seleção quantitativa. O avaliador, congelado no commit \texttt{34f28c5}, calculou o AUROC de imagem por postos e aproximou AUROC de pixel e AUPRO com histogramas de 2.048 bins. A AUPRO foi normalizada até FPR de 0,30. A média macro dá o mesmo peso a cada categoria.

\section{Análise dos Resultados}
\begin{table*}[htbp]
\centering
\caption{Resultados no teste oficial do MVTec AD em porcentagem.}
\label{tab:mvtec-final}
\setlength{\tabcolsep}{4pt}
\begin{tabular}{lrrrrrr}
\hline
 & \multicolumn{3}{c}{Baseline (época 109)} & \multicolumn{3}{c}{Completo (época 315)} \\
Categoria & AUROC img. & AUROC pix. & AUPRO & AUROC img. & AUROC pix. & AUPRO \\
\hline
bottle & 96,35 & 91,09 & 84,10 & 96,51 & 91,02 & 84,10 \\
cable & 71,68 & 73,33 & 65,75 & 71,01 & 73,12 & 65,88 \\
capsule & 64,62 & 93,04 & 84,32 & 65,10 & 93,12 & 84,56 \\
carpet & 89,73 & 94,88 & 86,28 & 89,17 & 94,50 & 85,81 \\
grid & 85,05 & 94,73 & 81,37 & 85,71 & 95,20 & 82,89 \\
hazelnut & 97,86 & 95,61 & 90,28 & 98,00 & 95,67 & 90,44 \\
leather & 94,53 & 97,51 & 93,95 & 94,16 & 97,51 & 94,02 \\
metal\_nut & 57,33 & 81,50 & 63,52 & 56,94 & 81,34 & 63,52 \\
pill & 66,18 & 86,07 & 86,22 & 66,83 & 86,19 & 86,35 \\
screw & 51,22 & 95,46 & 84,24 & 51,71 & 95,54 & 84,48 \\
tile & 94,08 & 87,19 & 68,45 & 93,80 & 86,82 & 67,95 \\
toothbrush & 80,00 & 92,35 & 72,28 & 80,56 & 92,59 & 72,87 \\
transistor & 56,29 & 65,99 & 43,38 & 56,17 & 66,06 & 43,77 \\
wood & 97,63 & 93,05 & 91,96 & 97,54 & 93,00 & 91,93 \\
zipper & 73,45 & 94,38 & 82,25 & 73,69 & 94,50 & 82,67 \\
\hline
Média macro & 78,40 & 89,08 & 78,56 & 78,46 & 89,08 & 78,75 \\
\hline
\end{tabular}
\end{table*}


O modelo completo ganhou 0,19 ponto percentual de AUPRO macro. Os dois AUROCs macro ficaram praticamente iguais. \texttt{screw}, \texttt{transistor} e \texttt{metal\_nut} tiveram os menores AUROCs de imagem do modelo completo. \texttt{transistor} teve também o menor AUPRO, 43,77\%.

\begin{figure*}[t]
\centering
\includegraphics[width=0.94\textwidth]{../analysis/final_eval_315/casos_objetos.png}
\caption{Casos de objetos com imagem, máscara, mapa e sobreposição. A cor indica discrepância L2 e o contorno ciano marca o defeito anotado. As linhas mostram um verdadeiro positivo, um falso negativo e um falso positivo.}
\end{figure*}

\begin{figure*}[t]
\centering
\includegraphics[width=0.94\textwidth]{../analysis/final_eval_315/casos_texturas.png}
\caption{Casos de texturas. O exemplo de falso alarme usa o percentil 95 dos escores normais do grupo, pois o limiar global não gerou falsos positivos em texturas.}
\end{figure*}

\begin{figure*}[t]
\centering
\includegraphics[width=0.94\textwidth]{../analysis/final_eval_315/escores_normais.png}
\caption{Escores das imagens normais por categoria, com escala horizontal compartilhada.}
\end{figure*}

As medianas dos escores normais variaram de 21,11 em \texttt{leather} a 32,60 em \texttt{toothbrush}. O limiar global de diagnóstico, definido no percentil 95 das normais da rodada, concentrou os 24 falsos positivos em objetos e deixou 369 de 382 defeitos em texturas abaixo do limiar. Esses valores descrevem o teste e dependem do limiar usado.

\section{Conclusões e Discussões}
Os mapas destacaram os furos na madeira, responderam fracamente à alteração de cor no couro e ativaram na cabeça e na rosca de um parafuso normal. Pose e reflexos são hipóteses para essa sensibilidade e não foram isolados experimentalmente. A grade de características de $32\times32$ e a interpolação a $512\times512$ limitam a localização fina. O extrator foi pré-treinado no ImageNet, a validação de treino continha apenas imagens normais e os resultados vêm de uma única semente. A diferença de escala entre categorias é mais evidente do que o ganho de AUPRO da variante completa.

\bibliographystyle{IEEEtran}
\bibliography{referencias}

\end{document}
